%% RNAAS submission draft, AASTeX v7.0.1. Compile on Overleaf (the AAS template ships aastex701.cls).
%% Every number traces to a file in this repository; see singles/rnaas/numbers.md.
\documentclass{aastex701}

\begin{document}

\title{Revisiting Stale TESS Single-transit Community Candidates: Mostly Eclipsing Binaries, Two Candidates with a Second Transit}

\author[orcid=0009-0000-7524-4470]{Jon McBee}
\affiliation{Independent researcher, USA}
\email{jon.mcbee@gmail.com}

\begin{abstract}
Community TESS Objects of Interest (CTOIs) flagged as single transits are rarely revisited.
I re-examined 303 such stars against all TESS data through Sector 108, with a validated repeat
search, difference imaging and Gaia DR3 binarity diagnostics. I recover the catalogued event for 272 stars; the other 31 mostly have a wrong
epoch, a weaker dip than catalogued, or no dip. Of the 272, 79 show repeats: 21 are Gaia-flagged binaries,
18 are blends, 23 fit no single period, and 17 needed vetting by hand. An any-depth search shows that the
CTOI 147035057.01 event is the secondary eclipse of a 37.315\,d eccentric binary
with a 5.5\% primary. Two on-target candidates with two transits survive: TIC 159614322, where a
partial third transit leaves $P=89.965$ or $179.930$\,d, and TIC 229055790. For 193 candidates that
remain single, a median 52\% of log-period space (5--1500\,d) is now excluded.
\end{abstract}

\keywords{\uat{Exoplanet detection methods}{489} --- \uat{Transit photometry}{1709} --- \uat{Eclipsing binary stars}{444}}

\section{Sample and Methods}

Of 794 open CTOIs with no period or tagged ``single'', ``duo'' or ``mono'', I excluded 344 from
groups still following up their own single transits, identified by ExoFOP uploader: members of the
TESS Single Transit Planet Candidate working group (223), the NGTS duotransit program
\citep[55;][]{Hawthorn2024} and the search of \citet{Salinas2025} (66). Of the other 450, I removed
46 with missing parameters, 8 on confirmed-planet hosts, 4 on TOI hosts and 24 on catalogued TESS
EBs \citep{Prsa2022}, leaving 368 CTOIs on 359 stars. Of these, 32 already have a community period
and 22 have $T<7$ or catalogued depths outside 200--50{,}000\,ppm. I analyzed the other 305 stars
plus one pilot target (TIC 53847727) later found in the EB catalogue; three truncated TIC IDs
duplicate other CTOIs, leaving 303 unique stars.

I used SPOC 2-min, TESS-SPOC \citep{Caldwell2020} and QLP \citep{Huang2020} light curves, and my own
TESScut aperture photometry \citep{Brasseur2019} for Sectors $\geq100$ without pipeline products.
Against SPOC 2-min data for ten stars ($T=8.5$--14.1, eight known planet hosts), my 8--9 pixel
aperture collects 0.78--0.99 of the SAP flux for $T<11.5$, and more for fainter stars, where SPOC
uses 4--7 pixels. Depths match PDCSAP to 3\% for uncrowded hosts (about 30\% shallower for crowded
ones, as I apply no crowding correction), mid-times agree to 2\,min, and the scatter is 1.07 times
SPOC's. The main difference is flux steps of a few ppt between segments.

A catalogued event is present if a dip deeper than $0.5\,\delta$ lies at SNR $\geq5$ within
$\pm0.03$\,d of the catalogued epoch ($\delta$, $D$: catalogued depth and duration). The repeat
search slides a box of that depth and duration on a $D/4$ grid with local baselines; a detection
needs SNR $\geq8$, a depth of 0.5--2$\,\delta$, 80\% coverage and baseline on both sides. The
matching is deliberate: only a recurrence of the same event enters $P=\Delta T/n$. Periods of
5--1500\,d come from linear-ephemeris fits, and any period predicting an epoch on flat data is
rejected. Validated blind from the first transit alone on 26 stars with two transits and a known
period, including duotransit candidates of \citet{Hawthorn2024}, the method kept the true period in
26 and found it uniquely in 19, with no wrong unique answers.

A source is off-target if TESScut difference images place it $\geq6''$ ($\geq3\sigma$) from the Gaia
DR3 position \citep{Gaia2023}; ``firm'' needs two agreeing images. Binarity indicators are Gaia DR3
\texttt{non\_single\_star}, two-body orbits and RV variability \citep[$p\leq0.01$, renormalised GOF
$>4$;][]{Katz2023}.

To find what the matched search misses, I also searched all 272 stars for dips of any depth
(1--24\,h; SNR $\geq8$ with red noise, 80\% coverage, baseline on both sides, no edge or momentum
dump, and a second reduction). On inverted light curves it finds 0.085 false events per
star-sector, against 0.092 in the real data, so individual extra events need vetting by eye. It
re-detects 234 of the 272 catalogued events (most misses are below SNR 8) and the secondary eclipses
of six known TESS EBs.

\section{Results}

The event test used pipeline light curves (SPOC 2-min for 285 stars), never my photometry; 31 stars
fail it. Inspected in every light curve (SPOC SAP and PDCSAP, TESS-SPOC, QLP, six other reductions
and mine), none fails because of my photometry, though in three marginal cases another pipeline
product would pass. Nine CTOIs have wrong epochs (six are 0.5\,d early), 13 show only a dip below
$0.5\,\delta$ or SNR 5, 6 show no dip, and 3 have flagged data, a blended neighboring CTOI or no data
at the epoch; 18 of the 31 come from one uploader's batch.

Of the 272 stars with the catalogued event, 193 show no repeat. Of the 79 with repeats, 21 are
Gaia-flagged binaries (three Gaia SB1 orbits reproduce the TESS periods; Table~\ref{tab:systems}),
18 are firmly off-target, 23 fit no single period of 5--1500\,d, and 17 were vetted by hand: 2
survive (below), 3 remain tentative (seen only in my Sector~107 photometry), 5 are possibly
off-target and 7 were rejected.

Vetted extra eclipses from the any-depth search appear on 14 stars: nine binaries, four of the 23
without a single period, and one single star, TIC 250492521, whose differently shaped event makes it
an EB or blend candidate. No single star gains a period. The extra eclipses give periods for four
stars, including TIC 188835620 (22.217\,d) and TIC 83384234 (4.33\,d); for two, the catalogued event
is a secondary eclipse.

\textit{CTOI 147035057.01.} Three matching events (Sectors 7, 33 and 108, the last in my
photometry) give $P=37.31497\pm0.00003$\,d, but they are secondary eclipses: a V-shaped primary
5.3--5.6\% deep falls at phase 0.239 in Sectors 6 (TESS-SPOC) and 87 (QLP), 59 cycles apart. Gaia
sources within $60''$ supply only 3.5\% of the flux, so the deep eclipse is on target, and the Gaia DR3 RV
amplitude of 50\,km\,s$^{-1}$ makes the companion stellar. Grazing geometry in the eccentric orbit
($e\cos\omega\approx0.4$) could explain the depth ratio; Gaia DR4 epoch RVs should decide.

\textit{TIC 159614322.} Transits in Sectors 14 and 41 ($\approx3.2$\,ppt, $\approx9.7$\,h) are
consistent across products and on target, leaving $P=719.72/n$ with $n\in\{1,2,4,5,8\}$. Sector 21
shows an ingress and $\approx9$\,h of transit floor before the orbit-1 downlink gap, in TESS-SPOC,
QLP and my photometry, with a flat background and no pointing step; a blind re-analysis with
independent code also found it. It selects $n=8$ or 4: $P=89.9651\pm0.0006$ or
$179.9302\pm0.0012$\,d, $T_0={\rm BJD_{TDB}}\,2458703.8277\pm0.0041$ (fitted to the full transits;
the partial one is 13--26\,min late, perhaps biased by the missing egress). The companion is
$\approx0.9\,R_{\rm J}$ and the host is RV-stable in Gaia DR3. Both periods predict a transit at
2026-12-28 06:40\,UTC ($\pm45$\,min, errors inflated $\times2.1$).

\textit{TIC 229055790.} This Planet Hunters TESS candidate \citep[see][]{Eisner2021} has consistent,
on-target transits in Sectors 1 and 68, re-detected independently. With 38 allowed periods
$\geq26.7$\,d and $R_p\approx3.8\,R_\oplus$, it is worth monitoring.

\begin{deluxetable*}{llll}
\tablecaption{Selected systems.\label{tab:systems}}
\tablehead{\colhead{TIC} & \colhead{Events (Sectors)} & \colhead{Period} & \colhead{Notes}}
\startdata
147035057 & 7, 33, 108$^a$; 6, 87$^c$ & $37.31497\pm0.00003$\,d & EB; Gaia RV-variable \\
159614322 & 14, 41; 21 (partial) & $89.9651\pm0.0006$ or $179.9302\pm0.0012$\,d & $\approx0.9\,R_{\rm J}$; Gaia RV-stable \\
229055790 & 1, 68 & $1838.945/n\geq26.7$\,d (38 aliases) & $\approx3.8\,R_\oplus$ \\
231969683 & 6, 32, 34, 87, 98 & 21.978\,d & Gaia SB1: $21.969\pm0.005$\,d \\
230130556 & 32 sectors, 15--86 & 22.259\,d & Gaia SB1: $22.245\pm0.024$\,d \\
458419988 & 48, 75 & $734.56/40=18.36$\,d & Gaia SB1: 18.355\,d \\
268534931 & 3, 30; 107$^{a,b}$ & 92.60, 46.30, 30.87 or 23.15\,d$^b$ & \citet{Hawthorn2024} candidate \\
139726266 & 32; 107$^{a,b}$ & $2109.0/n$\,d$^b$ & $\approx2.8\,R_\oplus$ \\
\enddata
\tablecomments{$^a$My TESScut photometry. $^b$Tentative until Sector 107 pipeline products appear.
$^c$Primary eclipse (5.3--5.6\%) at phase 0.239; the catalogued event is the secondary.}
\end{deluxetable*}

\section{Conclusion}

Stale single-transit CTOIs with repeats are mostly stellar; Gaia binarity checks and difference
imaging dispose of half of them. Two on-target candidates with short alias lists remain, plus
tentative repeats that pipeline products will settle. Code and classifications of all 306 CTOIs are
at \url{https://github.com/coracle-us/duo-transits}.

\begin{acknowledgments}
This work uses TESS data from MAST: SPOC \dataset[doi:10.17909/t9-nmc8-f686]{https://doi.org/10.17909/t9-nmc8-f686},
TESS-SPOC \dataset[doi:10.17909/t9-wpz1-8s54]{https://doi.org/10.17909/t9-wpz1-8s54} and QLP
\dataset[doi:10.17909/t9-r086-e880]{https://doi.org/10.17909/t9-r086-e880} light curves and
full-frame images \dataset[doi:10.17909/0cp4-2j79]{https://doi.org/10.17909/0cp4-2j79}; Gaia DR3;
ExoFOP \dataset[doi:10.26134/ExoFOP5]{https://doi.org/10.26134/ExoFOP5}; and the NASA Exoplanet
Archive \dataset[doi:10.26133/NEA13]{https://doi.org/10.26133/NEA13}.
\end{acknowledgments}

\facilities{TESS, Gaia}
\software{lightkurve \citep{Lightkurve2018}, astroquery \citep{Ginsburg2019}, batman \citep{Kreidberg2015}, astrocut \citep{Brasseur2019}}

\section*{AI disclosure statement}
I used Anthropic's Claude, an AI assistant, under my supervision to write and run the analysis code
and to draft and edit this note. I reviewed the text, checked every number against the saved
outputs and every reference against ADS and arXiv, and take full responsibility for the content.
The note contains no AI-generated figures.

\begin{thebibliography}{}
\bibitem[Brasseur et al.(2019)]{Brasseur2019} Brasseur, C.~E., Phillip, C., Fleming, S.~W., Mullally, S.~E., \& White, R.~L. 2019, Astrocut: Tools for creating cutouts of TESS images, Astrophysics Source Code Library, ascl:1905.007
\bibitem[Caldwell et al.(2020)]{Caldwell2020} Caldwell, D.~A., Tenenbaum, P., Twicken, J.~D., et al. 2020, RNAAS, 4, 201
\bibitem[Eisner et al.(2021)]{Eisner2021} Eisner, N.~L., Barrag\'an, O., Lintott, C., et al. 2021, MNRAS, 501, 4669
\bibitem[Gaia Collaboration et al.(2023)]{Gaia2023} Gaia Collaboration, Vallenari, A., Brown, A.~G.~A., et al. 2023, A\&A, 674, A1
\bibitem[Ginsburg et al.(2019)]{Ginsburg2019} Ginsburg, A., Sip{\H o}cz, B.~M., Brasseur, C.~E., et al. 2019, AJ, 157, 98
\bibitem[Hawthorn et al.(2024)]{Hawthorn2024} Hawthorn, F., Gill, S., Bayliss, D., et al. 2024, MNRAS, 528, 1841
\bibitem[Huang et al.(2020)]{Huang2020} Huang, C.~X., Vanderburg, A., P\'al, A., et al. 2020, RNAAS, 4, 204
\bibitem[Katz et al.(2023)]{Katz2023} Katz, D., Sartoretti, P., Guerrier, A., et al. 2023, A\&A, 674, A5
\bibitem[Kreidberg(2015)]{Kreidberg2015} Kreidberg, L. 2015, PASP, 127, 1161
\bibitem[Lightkurve Collaboration et al.(2018)]{Lightkurve2018} Lightkurve Collaboration, Cardoso, J.~V.~d.~M., Hedges, C., et al. 2018, Lightkurve: Kepler and TESS time series analysis in Python, Astrophysics Source Code Library, ascl:1812.013
\bibitem[Pr{\v s}a et al.(2022)]{Prsa2022} Pr{\v s}a, A., Kochoska, A., Conroy, K.~E., et al. 2022, ApJS, 258, 16
\bibitem[Salinas et al.(2025)]{Salinas2025} Salinas, H., Brahm, R., Olmschenk, G., et al. 2025, MNRAS, 538, 2031
\end{thebibliography}

\end{document}
